\documentclass{article}

\usepackage[
  paperheight=8.5in,
  paperwidth=5.5in,
  left=10mm,
  right=10mm,
  top=20mm,
  bottom=20mm]{geometry}
\usepackage[utf8]{inputenc}

\usepackage{graphicx}
\usepackage{wrapfig}
\usepackage[bottom]{footmisc}
\usepackage{listings}
\usepackage{enumitem}

\usepackage{wrapfig}
\usepackage{ragged2e}

\usepackage{array}
\usepackage[table]{xcolor}
\usepackage{multirow}
\usepackage{booktabs}
\usepackage{hhline}
\definecolor{palegreen}{rgb}{0.6,0.98,0.6}

\usepackage{amsmath}
\usepackage{amssymb}
\usepackage{multicol}
\usepackage{lipsum}
\usepackage{hyphenat}
\PassOptionsToPackage{hyphens}{url}
\usepackage{url}

\usepackage{rotating}

%\usepackage{xeCJK}

%% support use of straight quotes in code listings
\usepackage[T1]{fontenc}
\usepackage{textcomp}
\usepackage{listings}
\lstset{upquote=true}

%% for shrinking space between lines
\usepackage{setspace}

\newcommand*{\affaddr}[1]{#1} % No op here. Customize it for different styles.
\newcommand*{\affmark}[1][*]{\textsuperscript{#1}}
\newcommand*{\email}[1]{\small{\texttt{#1}}}
\newcommand{\tarot}{\textsc{Tarot}}
\renewcommand*\contentsname{\centering Table of Contents}

\renewcommand{\footnoterule}{%
  \kern -3pt
  \hrule width \textwidth height 0.5pt
  \kern 2pt
}

% remove date
\date{}

\usepackage{titlesec}
\titleformat*{\section}{\large\bfseries}
\titleformat*{\subsection}{\normalsize\bfseries}
\titleformat*{\subsubsection}{\normalsize\bfseries}


\usepackage[backend=biber,style=numeric,sorting=none]{biblatex}
\addbibresource{quantum_tutorial.bib}

\title{A Study Of the Concepts, Capabilities, and Limits of Quantum Computing\\
\vspace{0.2in}
\large Conference Tutorial}

\author{
  Adam Wade Lewis\\
  Division of Mathematical, Computer, and Applied Sciences\\
  Athens State University\\
  Athens, Alabama 35611\\
  \texttt{Adam.Lewis@athens.edu}
}

\date{}

\begin{document}
\maketitle

Quantum computing is a model of computation that uses principles from
quantum mechanics to process information \cite{kaye_laflamme_mosca_2007}.
Classical computation begins with the bit, which stores a value of either
$0$ or $1$. Quantum computers instead use \textit{qubits}, whose states can
exhibit superposition, entanglement, and interference. These properties make
quantum computers useful for certain specialized problems, including integer
factorization and structured search, where quantum algorithms can offer
advantages over known classical approaches \cite{nielsen_chuang_2010}.

This hands-on tutorial offers a practical introduction to quantum computing
for undergraduate computing students, technically curious faculty, and IT
professionals. No prior background in quantum mechanics is required although
familiarity with basic computer science, linear algebra, and
probability will be helpful.

Participants will use a Jupyter Notebook \cite{Kluyver2016aa} built around
IBM's Qiskit framework \cite{qiskit_docs_2026,javadi_abhari_qiskit_2024}.
The notebook uses Qiskit to construct small quantum circuits, run them on
simulators, and interpret the resulting measurement counts. These examples
allow participants to see how quantum programs are represented, executed, and
analyzed without requiring access to physical quantum hardware.

The session begins by comparing the classical bit with the qubit, the basic
unit of quantum information. Participants then explore measurement,
superposition, probability amplitudes, and entanglement through small circuit
examples. The presentation avoids unnecessary physics and instead frames
quantum behavior in terms of information representation and computation.

The tutorial next introduces quantum gates and circuits as counterparts to
classical logic gates and circuits. Examples using the Hadamard gate and the
controlled-NOT gate show how quantum states are transformed and how simple
circuits can produce superposition and entanglement. The session then briefly
connects these ideas to landmark algorithms, including Grover's search
algorithm \cite{grover_1996} and Shor's factoring algorithm \cite{shor_1994}.

The final portion addresses the practical limits of current quantum
computing. Participants consider what quantum computers may do well, what
they do not yet do well, and why noise, decoherence, limited qubit counts,
and error correction remain central challenges. The discussion also connects
quantum computing to likely application areas, including cryptography,
optimization, quantum simulation, and quantum-safe security
\cite{nist_pqc_standards_2024}.

By the end of the tutorial, participants should be able to distinguish a
classical bit from a qubit, describe superposition and entanglement in plain
language, explain the role of quantum gates and circuits, identify several
important quantum algorithms, and separate realistic near-term uses of quantum
computing from common misconceptions.

\section*{Availability}
The Jupyter Notebook and supporting files for this tutorial are hosted on
GitHub at \url{https://github.com/adamwadelewis/quantum_tutorial}.

\section*{Biography}
Dr. Lewis joined Athens State University's Division of Mathematical,
Computer, and Applied Sciences in the Fall of 2012 after an extensive
career in the computing industry. His research interests include
computational linguistics, high-performance computing, and the computer
architectures and systems that support those fields.

\medskip

{\small
\printbibliography
}

\end{document}
